% 问题一、问题二摘要；可置于主文档的 abstract 环境内。
% 两段纯文本合计约 360 字（含数字、标点）；不含总述、问题三四及关键词。
% “线性时间”仅指有序凸多边形顶点给定后的直径求解阶段。
% 问题二的数值为名义方位近似下的离散采样评价值，不是实际定位误差上界。

\textbf{针对问题一，}建立考虑 $\pm1^\circ$ 测向误差的\textbf{半平面交会定位模型}，利用凸性将区域直径计算转化为最远顶点对搜索，并采用\textbf{旋转卡壳算法}在线性时间内求解。在此基础上，证明同直径覆盖圆的圆心只能为最远点对中点，给出逐顶点覆盖判据。构造三次测向形成的等边三角形反例，得到区域直径为 $20\,\mathrm{m}$、最小覆盖圆直径约 $23.09\,\mathrm{m}$，表明以区域直径为直径的圆\textbf{不一定能完全覆盖定位区域}。

\textbf{针对问题二，}结合首次示向度、接收距离及场地边界筛选备选测点，以源位置采样下的最大名义交会直径为指标，建立\textbf{极小极大选址模型}。通过网格搜索与角域边界采样求解，将评价值不超过网格最优值 $1.05$ 倍的测点组成近优候选区域。沿半径朝外、斜穿圆域与近切向斜穿三个算例的最优采样评价值分别为 $47.83$、$64.37$ 和 $22.67\,\mathrm{m}$；固定所选测点加密源采样后，评价值最大变化约 $0.41\%$。
